\documentclass[10pt,a4paper]{amsart}
\usepackage[margin=19mm]{geometry}
\usepackage{amsmath,amssymb,mathtools}
\usepackage[scr=rsfs,cal=euler]{mathalfa}
\usepackage{xfrac}
\usepackage[unicode,hidelinks,bookmarks=false]{hyperref}
\newcommand{\ie}{i.e.\ }
\newcommand{\cf}{cf.\ }
\newcommand{\resp}{respectively\ }
\newcommand{\Subsection}{\subsection}
\providecommand{\nobreakdash}{\nobreak\hbox{-}}
\setlength{\emergencystretch}{2em}
\multlinegap0pt
\newcommand{\PGL}{\operatorname{PGL}}
\usepackage{amssymb}
\usepackage{mathtools}
\usepackage{xcolor}
\usepackage[shortlabels]{enumitem}
\newtheorem{lemma}{Lemma}
\newtheorem{theorem}{Theorem}
\newcommand{\RR}{\mathbb R}
\newcommand{\ZZ}{\mathbb Z}
\title{Vertex volumes, lattice-minima tails, and height zeta functions for the standard arithmetic quotient of $\PGL_d$}
\author{Soonki Hong, Sanghoon Kwon}
\date{Source: arXiv:2607.28433v1 (CC BY 4.0)}
\begin{document}
\maketitle

\noindent\emph{Source and licence.} Vertex volumes, lattice-minima tails, and height zeta functions for the standard arithmetic quotient of $\PGL_d$. Version arXiv:2607.28433v1 (CC BY 4.0), distributed by arXiv under the Creative Commons Attribution 4.0 International licence, http://creativecommons.org/licenses/by/4.0/, as recorded in the arXiv licence metadata for this exact version and shown on the abstract page https://arxiv.org/abs/2607.28433v1. Full licence notice: \texttt{licenses/arxiv-2607.28433-CC-BY-4.0.txt}.\par\smallskip

\noindent\textit{Editorial scope.} Counted unit: the article's theorem thm:sharp-tail (source0.tex lines 871--879). The article's complete original proof of it is delivered with it (source0.tex lines 881--929, one \texttt{proof} environment of 49 lines). No mathematics is written, repaired or re-derived: the statement and the proof are the article's own lines, in the article's own notation, and no retained line is substituted or re-flowed.  Retained uncounted as the source-local context the proof needs: lines 786--796.  Nothing was omitted from any retained span, so the package carries no omission record.  No retained line contains a citation, so the wrapper carries no bibliography and no bibliography entry.  No retained line contains a figure environment, an image-inclusion command or drawing code, so no image is delivered and the package declares no asset dependency.  Editorial formatting only, in addition: an AMS \texttt{amsart} preamble that restates the article's own declarations and macros used by the retained lines, namely the environments \texttt{lemma}, \texttt{theorem} on separate counters, together with the standard packages those lines need.  The retained environments print in this document as Lemma 1, Theorem 1 in order of appearance, whereas the article prints the same items with the numbering of its own full document.  The complete native source of the article as distributed in the arXiv e-print for this exact version is bundled unchanged as \texttt{sources/206418/source0.tex}.
\begin{lemma}[Transverse gap]\label{lem:transverse-gap}
Fix $1\le j\le d-1$.  For every $\delta\in\RR_{\ge0}^{d-1}$,
\begin{equation}\label{eq:Rj-formula}
  \Phi(\delta)-dS_j(\delta)
  =
  \sum_{i<j}i(j-i)\delta_i+
  \sum_{i>j}(i-j)(d-i)\delta_i.
\end{equation}
In particular this linear form is nonnegative and vanishes exactly on the ray
$\RR_{\ge0}e_j$.
\end{lemma}

\begin{theorem}[Sharp tail exponent]\label{thm:sharp-tail}
There exist constants $0<A_1\le A_2<\infty$, depending only on $d$ and $q$, such that
\begin{equation}\label{eq:sharp-tail}
  A_1T^{-d}\le \nu(E_T)\le A_2T^{-d}
  \qquad (T\ge1).
\end{equation}
The same estimate holds for the probability measure $\mu$ after multiplying the constants
by $\nu(Y)^{-1}$.
\end{theorem}

\begin{proof}
Let $C_*=\max_I C(I)$, where $I$ runs over all subsets of $\{1,\ldots,d-1\}$.  The upper
bound follows by decomposing according to a partial sum that realizes $H$.  For each
$\delta$ with $H(\delta)>L$, choose one index $j$ such that $H(\delta)=S_j(\delta)$.  Then
\[
  \nu(E_T)
  \le
  C_*\sum_{j=1}^{d-1}
  \sum_{\substack{\delta\in\ZZ_{\ge0}^{d-1}\\ S_j(\delta)>L}}
  q^{-\Phi(\delta)}.
\]
By Lemma~\ref{lem:transverse-gap}, write
\[
  \Phi(\delta)=dS_j(\delta)+R_j(\delta),
\]
where
\[
  R_j(\delta)=\sum_{i<j}i(j-i)\delta_i+
  \sum_{i>j}(i-j)(d-i)\delta_i.
\]
The form $R_j$ does not involve $\delta_j$ and is strictly positive in every transverse
coordinate.

Fix the transverse vector $z=(\delta_i)_{i\ne j}$.  Let
\[
  M_j(z)=\sum_{i\ne j}b_{ji}z_i,
  \qquad c_j=j(d-j).
\]
Since $dS_j(\delta)=c_j\delta_j+M_j(z)$, the sum over all $\delta_j\ge0$ satisfying
$S_j(\delta)>L$ is bounded by
\[
  \sum_{\delta_j:\,c_j\delta_j+M_j(z)>dL}
  q^{-c_j\delta_j-M_j(z)}q^{-R_j(z)}
  \le C_jq^{-dL}q^{-R_j(z)}
\]
for a constant $C_j$ depending only on $d$.  Summing over all transverse $z$ gives a finite
geometric product because all nonzero coefficients of $R_j$ are positive.  Hence
$\nu(E_T)\ll q^{-dL}=T^{-d}$.

For the lower bound, use the ray $\delta=Ne_1$.  Choose
\[
  N_T=\left\lfloor\frac{dL}{d-1}\right\rfloor+1.
\]
Then $H(N_Te_1)>L$, so this single vertex lies in $E_T$.  Its measure is
\[
  C(\{1\})q^{-(d-1)N_T}\gg q^{-dL}=T^{-d}.
\]
This proves \eqref{eq:sharp-tail}.
\end{proof}

\end{document}
